% ==============================================================================
% 06_fazit_ausblick.tex – Kapitel 6
% ==============================================================================
\chapter{Fazit und Ausblick}
\label{ch:fazit}

\section{Fazit}
\label{sec:fazit}
Diese Arbeit hat untersucht, wie gut man ein vortrainiertes
Bildverarbeitungsmodell mit einer kurzen Webcam-Kalibrierung an eine Person
anpassen kann. Ein auf ImageNet vortrainiertes ResNet-18 blieb eingefroren,
und nur ein kleiner Regression Head wurde trainiert.

Die wichtigsten Ergebnisse sind:
\begin{itemize}
  \item Ein allgemeines Modell, das auf GazeCapture trainiert wurde, ist nicht
        besser als eine Schätzung, die immer die Bildschirmmitte ausgibt. Es
        dient daher vor allem als Startpunkt für die Personalisierung.
  \item Mit einem realistischen, zeitlichen Split senken wenige Minuten
        persönlicher Daten den Fehler von \SI{31,8}{\percent} auf
        \SI{27,9}{\percent} der Bildschirmdiagonale (etwa $-$\SI{12}{\percent},
        18 Sitzungen). In guten Sitzungen bleibt der Gewinn groß (etwa
        $-$\SI{46}{\percent}) und gilt auch für eine neue Aufnahme derselben
        Person. Ein Lauf dauert etwa zwei Sekunden auf einer CPU.
  \item Ein zufälliger Split ergibt \SI{19,8}{\percent} statt
        \SI{30,3}{\percent} (etwa $-$\SI{35}{\percent}). Er überschätzt den
        Nutzen aber, weil fast gleiche Nachbarbilder in Training und Test
        liegen.
  \item Ruhiger Hintergrund und gutes Licht sind wichtiger als die Gestaltung
        der Kalibrierung. Eine gleichmäßige Bewegung des Ziels über den ganzen
        Bildschirm war am besten.
  \item Bilder direkt nach einem Sprung des Ziels haben falsche Labels. Sie zu
        entfernen hilft wenig, kostet aber nichts.
\end{itemize}
Insgesamt zeigt die Arbeit: Dateneffiziente Personalisierung funktioniert mit
einfachen Mitteln. Die Genauigkeit hängt aber stark von den
Aufnahmebedingungen ab, und man muss sie mit einer realistischen Aufteilung der
Daten bewerten. Alle Aufnahmen stammen von einer Person (dem Autor). Die
Ergebnisse gelten daher zunächst nur für diese Person; ob sie sich auf andere
Personen übertragen lassen, muss eine Studie mit mehreren Personen zeigen.

\section{Ausblick}
\label{sec:ausblick}
Mehrere nächste Schritte liegen nahe:
\begin{itemize}
  \item \textbf{Augen- oder Gesichtsausschnitte als Eingabe.} Im ganzen
        Kamerabild sind die Augen sehr klein. Ein Ausschnitt sollte dem Modell
        viel mehr Blickinformation geben.
  \item \textbf{Andere vortrainierte Modelle.} DINOv2, CLIP, MAE oder UniGaze
        statt ResNet-18, jeweils mit eingefrorenem Backbone und kleinem Head.
  \item \textbf{Sparsame Anpassung des Backbones.} Adapter oder LoRA als
        Mittelweg zwischen reinem Head-Training und vollständigem Fine-Tuning.
  \item \textbf{Mehr Personen und mehr Sitzungen,} mit Bedingungen, die sich
        jeweils nur in einem Faktor ändern.
  \item \textbf{Kopfbewegung berücksichtigen,} zum Beispiel mit der
        Kopfposition als zusätzlichem Feature oder mit normalisierten
        Gesichtsbildern.
\end{itemize}

\paragraph{Web-Demo.}
Als praktischen Test habe ich eine Demo im Browser gebaut (Server mit FastAPI,
geeignet für Hugging Face Spaces). Der Nutzer folgt 60 Sekunden lang einem
Punkt. Danach wird ein Modell personalisiert, und der Blick wird live als
Lichtfeld auf einem Raster von $3\times3$ Feldern angezeigt. Am Ende zeigt eine
Auswertung, wann welches Feld angeschaut wurde. Die Demo nutzt bereits einige
Ideen aus dem Ausblick: einen Augenausschnitt, räumliche Features des
ResNet-18, eine Ridge Regression und die Kopfposition als Zusatz-Feature. Ein
kurzer Test mit fünf Punkten (je \SI{3}{\second}) nach der Kalibrierung misst
den Fehler auf neuen Bildern. In mehreren Einzeltests lag dieser Fehler
zwischen etwa \SI{17,6}{\percent} und \SI{25,2}{\percent} (letzter Test:
\SI{19,0}{\percent}); eine konstante Schätzung erreicht auf den
Kalibrierdaten etwa \SI{23}{\percent}. In einem direkten Vergleich sank der
Testfehler mit der Kopfposition von \SI{23,0}{\percent} auf
\SI{17,6}{\percent}. Die Werte schwanken stark zwischen den Durchläufen und
wurden nicht systematisch ausgewertet. Sie zeigen aber, dass Kopfbewegung und
Licht im Live-Betrieb große Fehlerquellen sind.
